\subsubsection{Cliente}

Como se habló de las capas de la arquitectura del sistema del juego, éste se divide principalmente en dos: 
cliente y servidor, cada una con su propia lógica y responsabilidades. El cliente es responsable de la 
experiencia del usuario, incluyendo la representación visual, la interacción y la comunicación con el 
servidor. El servidor, por otro lado, maneja la lógica del juego, la gestión de datos y la coordinación 
entre múltiples clientes.

En esta sección se describirá la implementación del cliente del juego, detallando su estructura, componentes 
y cómo se integra con el resto del sistema. Se explicará la organización del proyecto en \texttt{Unity}, la 
arquitectura de sus módulos, la comunicación con el servidor y las herramientas utilizadas para su desarrollo.

El cliente contiene tres módulos principales: el \texttt{Core}, que contiene la lógica central y utilidades; 
el \texttt{Gameplay}, que maneja la lógica de juego en tiempo de ejecución; y el \texttt{User Interface}, 
que se encarga de la representación visual y la interacción del usuario. Para más detalle sobre modulos, 
\textit{Assemblies} y \textit{Vcontainer}, esto se explica en la sección 
(TODO: (reference) 8.2 Estructura del Proyecto de Unity). 

\paragraph{Cliente: Core}

El primer módulo base, el \texttt{Core} del cliente, está conformado principalmente por clases que perduran 
en memoria durante la escena y pueden ser utilizados por múltiples recursos.

Dentro de estos hay elementos de como datos/valores fijos ya inicialmente configurados para el correcto 
funcionamiento. Por otro lado también otros recursos importantes son los eventos que utiliza el cliente. 
Todos los eventos formalizan un \textit{event channel}, que es un canal de comunicación entre distintos 
componentes del cliente, donde se define para comunicar cambios entre los componentes, invocando y 
suscribiendose a estos eventos.

Todos estos elementos se encuentran inyectados en el contenedor de dependencias de \textit{VContainer}, 
lo que permite una gestión eficiente de las dependencias, una mayor modularidad y no depender de \texttt{Unity}
para anexar componentes a los objetos de la escena.

De esta forma se permite que el código se mantenga desacoplado lo mayormente posible. Hay clases que no 
tienen conocimiento de otras ni a quien envían sus eventos, en los cuales se comparte o no lógica de 
alguna característica, dejando el código lo más limpio posible y evitando dependencias ante alteraciones 
de distintos componentes. 

\vspace{0.20cm}

Otro elemento primordial en la capa de \texttt{Core} es la implementación del sistema de inputs del cliente, 
el cual se basa en el sistema de inputs de \texttt{Unity}, pero con una capa de abstracción que permite una mayor 
flexibilidad y modularidad. Este sistema de inputs se integra con el resto del cliente a través de eventos, 
lo que permite que cualquier componente pueda reaccionar a las acciones del usuario sin necesidad de conocer 
la fuente de esos inputs.

Se encuentran en el \texttt{Core} también varios gestores de recursos como gestor de Menú, de controlador del 
\textit{mouse}, de cámara y otros, que se encargan de manejar aspectos específicos del juego y proporcionar una 
interfaz para interactuar con ellos.

Dentro del módulo de \texttt{Core}, también hay un \textit{assembly} llamado \texttt{Common}, donde se guarda 
la lógica y los datos que deben ser idénticos para cliente y servidor. 

En dicho módulo, no contiene código de \texttt{UI} ni de \texttt{Gameplay} exclusivo, sino contratos, 
configuraciones y mecanismos de red que ambos lados usan para comunicarse de forma coherente. 
Esto garantiza que interpreten igual los valores de configuración y protocolos. 
Entre estos se encuentran \textit{DTO's}, \textit{Enums}, eventos, y otros elementos que son compartidos. 

\paragraph{Cliente: Gameplay}

El módulo \texttt{Gameplay} es la capa de jugabilidad y ejecución específica del cliente. No es solo lógica de red ni 
solo \texttt{UI}, es el conjunto que une la experiencia del jugador con los datos de mundo, con el servidor y 
con los objetos lógicos y visuales que se representan lo que el usuario percibe.

Todo comienza con el código en sus puntos de entrada, llamados \textit{Entry Points}, donde se maneja el arranque 
del cliente y toda la configuración necesaria para que la escena tome todo lo necesario para comenzar a funcionar y
tomar los recursos para que se mantenga así. 

Esto incluye la configuración de iniciar la conexión con el servidor, cargar los recursos necesarios ya sea local 
o desde el \textit{backend} que contiene los datos a cargar o renderizar, configurar los sistemas de juego y 
preparar la escena y su flujo para la interacción del usuario.

Todo estos recursos nombrados, previamente se inicializan e inyectan los recursos y servicios necesarios 
en el contenedor de dependencias de \textit{VContainer}, bajo clases inicializadoras de cada escena.
Con respecto a lo nombrado sobre escenas, en el cliente se manejan distintas escenas para cada parte del juego, 
una escena para el menú principal, que incluye el navegador de mundos, edición de personaje, menú de obtención 
de moneda de juego, etc. Y luego otra escena para el juego en sí, donde se desarrolla la experiencia de juego 
una vez que el usuario ha seleccionado un mundo y ha comenzado a jugar.

\vspace{0.25cm}

Volviendo a la configuración de la escena del juego, Cada objeto/entidad carga sus datos por su propio sistema 
categorizado de \textit{loaders}. Una vez que se ha establecido la conexión con el servidor y 
se han cargado los recursos necesarios, el \textit{Entry Point} del juego se encarga de configurar todos los 
objetos presentes en el juego mediante \textit{linkers}.
Estos son clases que se encargan de vincular al objeto/entidad del mundo con su representación visual y lógica 
en la escena.

Es una capa de \textit{“composition wiring”}: toma objetos de red y los convierte en entidades 
jugables con su \textit{UI} y comportamiento. Esto permite inyectar dependencias y controllers en cada tipo 
de entidad, diferenciar entre jugador local y otros clientes, y dejar las entidades listas para que el sistema 
de juego las maneje sin necesidad de acoplar la lógica de red o de \textit{UI} directamente en ellas, respondiendo 
a lo que el servidor envía y a lo que el jugador hace.

Dentro de este \textit{liking} un elemento importante, si bien declarado en el módulo de \texttt{Core}, se 
agrega a cada entidad el componente \textit{network adapter}, que actúa como el puente directo entre esa 
entidad y la red.
Para el jugador local, permite que esa entidad envíe acciones y transacciones al servidor. Para entidades 
remotas, permite recibir eventos del servidor y procesarlos localmente.
De esta forma así llega información como ataques, diálogos, portales, interacciones, etc. y se procesan en 
cada entidad para que el jugador vea la respuesta a sus acciones o a las acciones de otros jugadores.

Con esto se logra manejar el comportamiento de cada entidad desde la lógica en el servidor, explicado en su sección, 
teniendo el control según su \textit{NetId} almacenada en su \textit{network adapter} de cada entidad.
De esta forma el cliente solo se encarga de representar y ejecutar lo que el servidor le indica.

Para el módulo \texttt{Common} dentro de \texttt{Gameplay}, los \textit{State Storage} que contiene estados 
compartidos para personajes, ítems, inventario, etc. Estos componentes guardan el estado replicado o 
usado en red y son comunes a ambos servidor y cliente. Por ejemplo, el estado de salud de un personaje, su 
posición, o el contenido de un inventario.

Además en este módulo en común, se encuentran los objetos llamados \textit{Prefabs} que son objetos 
instanciables en la escena, que se comparten para tener una estructura similar y componentes basicos de red 
entre el cliente y servidor, y luego cada parte configura correctamente sus componentes o sistemas.

\paragraph{Cliente: User Interface}

El módulo de \textit{UI} del cliente es la capa que entrega toda la experiencia visual e interactiva del jugador 
fuera de la simulación de jugabilidad. Su función es presentar información, opciones y controles al 
usuario, y traducir las acciones de la interfaz en eventos que el resto del cliente consume.

Esto incluye elementos que visualizan el estado del juego, como barras de salud, inventario, 
menús de misiones, diálogos, etc. También está involucrado todo el flujo de navegación como el menú principal, 
selección de mundo, y el contenido gráfico interactuable o no dentro de la escena.

Se integra con el resto del cliente mediante eventos y controladores de cada elemento de \textit{UI}. De esta forma, 
al tener la \textit{UI} en su propio módulo, el cliente puede cambiar la presentación sin tocar la lógica de 
\textit{gameplay} o red, por lo que no hay dependencias entre la lógica de juego en el módulo \texttt{Gameplay} 
y la presentación visual.

